% DIAPOSITIVA 1: Contexto Histórico y Concepto Clave
\begin{frame}{Teorema de Szemerédi: Contexto Histórico}

    \begin{columns}

        \begin{column}{0.6\textwidth}
            Propuesto como conjetura por \textbf{Erdős y Turán} en
            1936, no fue hasta \textbf{1975} que \textbf{Endre
            Szemerédi} logró demostrarlo, obteniendo el
            \textbf{Premio Abel}.

            \vspace{0.4cm}

            \begin{block}{Concepto Clave: Densidad vs.\ Color}
                Mientras que Ramsey trabaja con \textit{coloraciones},
                Szemerédi trabaja con \textbf{densidad}: si un
                subconjunto de enteros es lo suficientemente
                ``grueso'', debe contener progresiones aritméticas.
            \end{block}
        \end{column}

        \begin{column}{0.4\textwidth}
            \begin{exampleblock}{Conexiones}
                \begin{itemize}
                    \item Teoría de números
                    \item Teoría ergódica
                    \item Informática (TCS)
                \end{itemize}
            \end{exampleblock}
        \end{column}

    \end{columns}

\end{frame}


% DIAPOSITIVA 2: Definición Formal
\begin{frame}{Teorema de Szemerédi: Enunciado Formal}

    \begin{alertblock}{Teorema (Szemerédi, 1975)}
        Sea $A \subseteq \{1, 2, \dots, N\}$. Para todo entero
        $k \geq 3$ y todo $\delta > 0$, existe $N(\delta, k)$ tal
        que si
        $$|A| \geq \delta N$$
        entonces $A$ contiene una \textbf{progresión aritmética de
        longitud $k$}.
    \end{alertblock}

    \pause
    \vspace{0.3cm}

    \begin{block}{Diferencia con Ramsey Tradicional}
        \begin{itemize}
            \item \textbf{Ramsey:} al dividir en $r$ colores, algún
                  color contendrá el patrón.
            \item \textbf{Szemerédi:} \textit{cualquier} conjunto con
                  densidad positiva contiene el patrón, sin necesidad
                  de considerar todas las clases.
        \end{itemize}
    \end{block}

\end{frame}


% DIAPOSITIVA 3: El Lema de Regularidad y Conclusión
\begin{frame}{El Lema de Regularidad de Szemerédi}

    Para demostrar el teorema, Szemerédi introdujo el
    \textbf{Lema de Regularidad}, hoy una herramienta fundamental
    en combinatoria:

    \begin{enumerate}
        \item \textbf{Partición en bloques:} Cualquier grafo grande
              puede dividirse en bloques más pequeños.
        \item \textbf{Comportamiento pseudo-aleatorio:} La mayoría de
              los bloques se comportan de manera casi aleatoria
              (``regular'').
        \item \textbf{Estructura forzada:} La regularidad hace
              inevitable la aparición de progresiones aritméticas.
    \end{enumerate}

    \pause
    \vspace{0.3cm}

    \begin{block}{Conclusión}
        El Teorema de Szemerédi refuerza la máxima de Ramsey:
        \textbf{la estructura es inevitable}. Si seleccionas una
        fracción constante de los enteros, no importa cuán caótica
        sea tu elección; terminarás creando una progresión aritmética
        perfecta, por \textit{pura necesidad matemática}.
    \end{block}

\end{frame}
